\chapter{Introduction}
\label{ch:introduction}

% 선행 연구 3편은 입문 본론 장으로 확장하지 않고, 본 논문의 문제의식이
% 발전한 과정만 간결하게 설명합니다.
\section{Research Goals and Significance}
\label{sec:ch1-goals}

Artificial-intelligence systems increasingly communicate trust through compact
claims.  A trained model may be described as differentially private, a generated
image may be labeled as authentic or watermarked, and a public interface may
return an acceptance decision.  Such statements are useful only when a reader can
determine what is protected or verified, which mechanism produced the relevant
result, and which evidence supports the wording.  A numerical value, a route
name, or a detector score is not by itself the guarantee that a public sentence
appears to express.

This distinction is especially important when a system crosses several semantic
layers.  In privacy-preserving learning, raw owners, events, and sequences may be
transformed into overlapping windows before a differentially private optimizer
operates on them.  The reported privacy parameters are meaningful only relative
to the privacy unit, neighboring relation, sampling rule, contribution bound, and
accounting procedure that generated them~\cite{dwork2014algorithmic,
kifer2020guidelines,ponomareva2023dpify,near2025guidelines}.  In generative-image
provenance, an observable watermark score may survive common transformations,
but the score does not alone establish which prompt, seed, model, or embedding
trace produced the image.  Signed content credentials add an important manifest
layer, yet their interpretation still depends on the technical relation actually
authenticated by the credential~\cite{longpre2024provenance,c2pa2026spec}.

The central goal of this dissertation is to reduce the gap between such public
claims and the system facts that make those claims defensible.  It introduces
\ClaimAlignedAssurance{} as a design perspective in which a positive trust claim
is released only when four elements remain mutually consistent: the intended
claim target, the executed or registered mechanism, the verification interface,
and the available evidence.  The framework is instantiated in two research axes:
\AuditablePrivacy{}, which concerns privacy claims about training, and
\SecurePublicVerification{}, which concerns independently verifiable provenance
for generative-AI content.

The first axis addresses privacy misalignment after a training run has completed.
A validator examines frozen evidence together with one requested privacy-unit
sentence and determines whether the evidence licenses the native accountant pair,
a registered converted consequence, or no positive wording.  This focus is
practical as well as semantic: a completed model may be reusable for a narrower
claim, may support a weaker converted statement, or may require retraining under a
mechanism native to the requested unit.  The decision cannot repair an
incompatible execution, but it can prevent the execution's isolated accountant
output from being silently promoted into a stronger public guarantee.

The second axis addresses a different trust object.  Public watermark detectors
enable independent screening, but the same public decision surface can be queried
or optimized by an adversary~\cite{jiang2023evading,muller2025blackbox,
lin2025crack}.  A private detector reduces this exposure at the cost of renewed
dependence on a central verifier.  This dissertation therefore studies a
conjunctive public decision: statistical image-side evidence is combined with a
cryptographic receipt bound to the declared generation context and committed
embedding trace.  The stronger acceptance statement is withheld unless both
conditions hold.

The significance of the dissertation lies in this alignment discipline rather
than in universal AI certification.  The methods do not prove correct execution,
universal privacy, or factual truth.  They make each positive statement narrower
and explicit about its authority.  Missing or inconsistent relations therefore
withhold positive wording rather than promote an incomplete signal.


\section{Research Background and Trajectory}
\label{sec:ch1-trajectory}

The present research problem developed through three preceding studies in
anonymous digital systems, decentralized authentication, and context-sensitive
AI generation.  These studies are not reproduced as independent technical
chapters in this dissertation.  Their role is to explain how the research focus
progressed from interpreting hidden structure, to reducing centralized trust, and
then to binding AI generation to domain context.  That trajectory motivates the
claim--mechanism--evidence question addressed in Chapters~\ref{ch:post-execution-validation}--\ref{ch:pp-mark}.

\subsection{Interpreting Hidden Roles in Anonymous Web3 Systems}
\label{subsec:ch1-trajectory-web3}

The first preceding study examined Ethereum, where addresses and transactions are
public but the service role of a smart contract and the relationship between
accounts are not directly stated.  It represented externally owned accounts,
contract accounts, and transactions as a heterogeneous graph and used a Graph
Attention Network version 2 model for DApp-category classification and account
link prediction~\cite{ko2024ethereum}.  The heterogeneous representation supplied
role context beyond an address's isolated balance and transaction values.

The antecedent for this dissertation is the separation between an observable
record and an asserted system role.  A graph model can infer a role from public
traces, but the inference is not an authenticated identity and remained bounded to
the evaluated Ethereum data.  The scope, model, and evidence determine what may
be said.  This gap between visible traces, hidden function, and warranted
interpretation became the first step toward claim-aligned reasoning.

\subsection{Reducing Centralized Trust in Metaverse Authentication}
\label{subsec:ch1-trajectory-metaverse}

The second preceding study moved from interpreting a hidden role to changing who
controls a verification decision.  Metaverse space access based on centrally
managed ranks or passwords leaves the service as both policy authority and point
of failure.  The proposed approach instead allowed a user to manage an
authentication token and used a blockchain smart contract and cosine similarity
to evaluate access to a declared virtual space~\cite{seo2024space}.  The study
calibrated the similarity threshold and tested feasibility in a metaverse setting.

The resulting lesson is that reducing centralized trust requires alignment among
the credential, spatial representation, policy, execution rule, and decision
context.  A smart contract can make a rule visible and repeatable, but blockchain
alone does not make every authentication statement trustworthy.  This work
therefore contributed a verification structure rather than a universal trust
claim.

\subsection{Preserving Domain Context in LLM Data Augmentation}
\label{subsec:ch1-trajectory-llm}

The third preceding study considered a data problem rather than access control.
Sensitive conversational domains can provide few authentic examples while
depending on evolving slang, implicit cues, and role-specific behavior.  Simple
text transformations may increase sample count without preserving this
interactional structure.  The proposed framework therefore encoded buyer and
seller personas, behavioral patterns, domain vocabulary, and dialogue rules in
prompts to a large language model~\cite{jeong2025persona}.  Multi-metric and expert
evaluation assessed whether the generated dialogues retained the intended domain
characteristics.

The antecedent is that an output-count claim is not a context-fidelity claim.  The
target and evidence must specify the roles, behavior, language rules, and
population being assessed.  The study remained limited to a Korean corpus
dominated by buyer--seller dyads, with evolving slang requiring updates.  A
favorable metric should not be detached from those conditions.

Taken together, the three studies trace a progression in the research question.
The Web3 study interpreted hidden roles from public traces; the metaverse study
reduced reliance on a centralized authentication decision; and the LLM study
bound generation to domain roles and linguistic behavior.  The present
dissertation advances that progression from inference, decentralization, and
context preservation to assurance.  It asks when a public privacy or provenance
statement is justified by the protected object, operative mechanism, verification
interface, and evidence that the system actually exposes.


\section{Claim-Aligned Assurance for Trustworthy AI}
\label{sec:ch1-caa}

\subsection{Claim--Mechanism--Evidence Alignment}
\label{subsec:ch1-alignment}

For this dissertation, a \emph{claim} is a positive, externally interpretable
statement about a defined object.  A privacy claim names a protected unit,
neighboring relation, and privacy parameters.  A provenance claim names an image,
generation context, and acceptance predicate.  A \emph{mechanism} is the set of
transformations, randomized procedures, samplers, executors, and cryptographic or
statistical operations from which the claim is derived.  A \emph{verification
interface} defines what a third party submits, observes, and interprets.  The
\emph{evidence} consists of the artifacts, commitments, runtime records,
accountant outputs, proofs, and consistency relations available to support that
interpretation.

Let (c), (M), (I), and (E) denote a requested claim, mechanism description,
verification interface, and evidence bundle, respectively.  Claim-Aligned
Assurance is represented as a scoped decision rather than a free-standing trust
score:

\begin{equation}
\mathsf{CAA}(c;M,I,E)
\in
\left\{
    \mathsf{Release}(c),
    \mathsf{Qualify}(c'),
    \mathsf{Block}
\right\},
\qquad c' \prec c.
\label{eq:ch1-caa-decision}
\end{equation}

The relation \(c'\prec c\) means that the evidence supports a narrower or weaker
statement than the one requested.  For example, frozen DP-SGD evidence may
support a generated-window statement but not an owner-level statement, or may
support a converted event-level parameter rather than the accountant's native
pair.  A watermark score may support statistical screening but not the stronger
statement that a declared generation trace has been publicly verified.  In both
cases, qualification preserves the distinction between what was observed and
what was established.

Table~\ref{tab:ch1-alignment-map} maps the common alignment elements to the two
research axes.  It is a dissertation-level synthesis, not an additional empirical
result.  The domains use different mechanisms and threat models, but each
positive statement must close a chain from its declared target to its released
decision.

\begin{table}[tbp]
    \centering
    \caption{Dissertation-wide alignment map for the two research axes}
    \label{tab:ch1-alignment-map}
    \footnotesize
    \setlength{\tabcolsep}{3.2pt}
    \renewcommand{\arraystretch}{1.14}
    \begin{tabular}{@{}
        >{\raggedright\arraybackslash}p{0.17\textwidth}
        >{\raggedright\arraybackslash}p{0.37\textwidth}
        >{\raggedright\arraybackslash}p{0.37\textwidth}@{}}
        \toprule
        \textbf{Alignment element} &
        \textbf{Auditable Privacy} &
        \textbf{Secure Public Verification} \\
        \midrule
        Claim target &
        Privacy unit, adjacency, and licensed
        \((\varepsilon,\delta)\) wording &
        Exact image, declared generation context, and provenance-acceptance
        statement \\

        Mechanism binding &
        Transformation, contribution rule, sampler, clipping/noise geometry,
        accountant, executor, and schedule &
        Context-derived payload, latent embedding, trace commitment, image-bound
        challenge, and proof predicate \\

        Verification interface &
        Post-execution wording query with an explicit unit and adjacency &
        Public score mode or conjunctive score-plus-receipt acceptance mode \\

        Evidence &
        Frozen artifacts, semantic digests, source/runtime relations, recomputed
        accounting, and registered obligations &
        Recoverable image signal, declared metadata, Merkle commitment, opened
        trace tuples, exact image hash, and SP1 receipt \\

        Fail-closed outcome &
        Block unsupported, unverified, invalid, or vacuous wording &
        Withhold full acceptance when either the score or receipt condition fails
        \\

        Principal boundary &
        Evidence-relative validation is not arbitrary-code proof or execution
        attestation &
        Accepted provenance evidence is not semantic truth, human authorship,
        legal identity, or proof of the complete diffusion execution \\
        \bottomrule
    \end{tabular}
\end{table}

Across these mappings, four design principles recur.  First, \emph{claim specification}
fixes the target, unit or context, predicate, and scope before a positive label is
interpreted.  Second, \emph{mechanism binding} connects that specification to the
actual route or generation process rather than to a label with a similar name.
Third, an \emph{evidence-bound decision} requires the relations needed by the
claim, rather than inferring the conclusion from an isolated numerical output.
Fourth, \emph{fail-closed assurance} suppresses a positive interpretation when a
required relation is absent, inconsistent, unregistered, or vacuous.

\clearpage

The framework treats the verification interface as part of the assurance object.
An interface determines which question is being answered and how its output may
be interpreted.  In the post-execution privacy study, the requested wording is an
explicit input because the same frozen run can support different answers for
window, event, and owner queries.  In PP-Mark, score mode and full acceptance are
separate interfaces because they establish different predicates.  Interface
separation prevents a weaker diagnostic from being silently presented as the
stronger claim.

\subsection{Assurance Boundaries and Fail-Closed Decisions}
\label{subsec:ch1-boundaries}

Claim alignment does not remove the need for a threat model.  Internally
consistent evidence may still come from a producer that did not execute the
reported code, and a proof establishes only the relation encoded by its circuit.
Every positive result is therefore relative to explicit premises, evidence
authority, and tested scope.

The privacy axis addresses an honest-but-fallible development setting.  Digests,
independent recomputation, and finite validation predicates expose registered
inconsistencies but do not attest hidden execution or arbitrary extensions.
The provenance axis considers adversarial signal creation, yet full acceptance
remains conditional on its hash, proof-system, key-control, statistical,
predicate, and partial-opening assumptions.  Its exact-image receipt must be
reissued after an edit, and its circuit neither re-executes complete diffusion nor
decides factual truth.

Fail-closed behavior marks an evidentiary boundary, not absence of the broader
property.  A blocked owner claim lacks licensing evidence, and a rejected PP-Mark
submission fails the declared score-and-receipt predicate.  Neither decision alone
proves that the broader privacy or provenance property is false.


\section{Research Questions and Contributions}
\label{sec:ch1-questions-contributions}

The dissertation addresses three research questions.

\begin{description}
    \item[\textbf{RQ1.}] Given a completed window-based DP-SGD execution and a
    requested privacy-unit sentence, which registered evidence obligations and
    transformations are sufficient to release, qualify, or block that exact
    wording?

    \item[\textbf{RQ2.}] How can generative-image provenance remain publicly
    verifiable while requiring more than an optimizable watermark score?

    \item[\textbf{RQ3.}] Which design principles are shared by these privacy and
    provenance settings, and where must their assurance conclusions explicitly
    stop?
\end{description}

\subsection{Auditable Privacy: Post-Execution Claim Validation}
\label{subsec:ch1-contrib-privacy}

Chapter~\ref{ch:post-execution-validation} addresses RQ1 after training.  Given a
frozen evidence bundle and wording query, it reopens the raw-to-generated mapping,
checks mechanism and accountant relations, and returns a direct, converted,
grouped, underspecified, or blocked path.  Its controlled conformance study, fixed
information-necessity witnesses, protocol-frozen activity-data executions, and
Sepsis and Backblaze mapping diagnostics test the decision rather than propose a
new accountant.  The resulting \AuditablePrivacy{} contribution is an inspectable
chain from a requested sentence to the observable premises that license or block
it.  It remains evidence-relative and is not promoted into end-to-end execution
attestation or a general proof of differential privacy.

\subsection{Secure Public Verification: Detection and Verifiable Provenance}
\label{subsec:ch1-contrib-provenance}

Chapter~\ref{ch:pp-mark} addresses RQ2 through PP-Mark.  A secret-keyed generation
binding determines the latent payload and sampled trace, which is Merkle-committed
and partially opened through an image-bound SP1 proof.  PP-Mark(score) remains a
statistical screen; PP-Mark(accept) requires both the calibrated score and a valid
receipt for the same context and exact image.  Forgery, removal, transformation,
cost, quality, and bounded model-generalization experiments evaluate the two
interfaces without equating either with semantic truth.

RQ3 is answered by the shared framework rather than by treating the two core
studies as one algorithm.  Across them, the dissertation contributes:

\begin{itemize}
    \item a common vocabulary for specifying a trust claim by its target,
    mechanism, verification interface, evidence, and assurance boundary;

    \item a post-execution, query-dependent validator that restricts privacy
    wording to what a frozen evidence chain licenses;

    \item a public provenance construction whose stronger decision joins a
    context-bound statistical signal with a publicly verifiable cryptographic
    receipt; and

    \item a cross-domain synthesis of claim specification, mechanism binding,
    evidence-bound decision, and fail-closed behavior, together with explicit
    non-claims for each instantiation.
\end{itemize}

The integrated contribution is therefore not a single privacy mechanism or a
single watermark detector.  It is a disciplined way to determine when a specific
trust statement remains aligned with the object, process, interface, and evidence
that give the statement meaning, demonstrated in two technically distinct AI
assurance domains.


\section{Dissertation Organization}
\label{sec:ch1-organization}

The remainder of the dissertation is organized as follows.

\paragraph{Chapter 2: Post-Execution Auditing of Privacy-Unit Claims.}
Chapter~\ref{ch:post-execution-validation} introduces the background of privacy
units and windowed DP-SGD, formalizes query-dependent evidence-bound validation,
and develops direct, converted, grouped, and blocked claim paths.  It evaluates
the contract through a 200-case controlled oracle, fixed necessity witnesses,
independently checked WISDM and UCI HAR executions, and Sepsis and Backblaze
mapping diagnostics.  The chapter also separates bundle consistency from
execution attestation and reports the exact authority of each experimental
surface.

\paragraph{Chapter 3: Secure Public Verification of Generative-AI Provenance.}
Chapter~\ref{ch:pp-mark} turns to the dissertation's second axis.  It describes
context-bound latent watermarking, deterministic trace commitment, SP1 proof
generation, and the separate score and full-acceptance interfaces of PP-Mark.  It
then evaluates forgery, removal, regeneration, transformation, efficiency,
quality, and bounded model-generalization conditions.  The chapter states the
cryptographic and statistical premises under which its public provenance
conclusion is interpreted.

\paragraph{Chapter 4: Conclusion.}
Chapter~\ref{ch:conclusion} summarizes the two core studies without repeating
their independent conclusions.  It relates evidence-bound privacy auditing to
secure public provenance verification and states the dissertation-wide design
principles of Claim-Aligned Assurance.  It concludes with limitations, broader
impact, and two concrete future validations: comparison of post-hoc conversion
with native patient-level DP, and retrieval of PP-Mark proof artifacts through a
C2PA-compatible public manifest repository.

This organization preserves the technical independence of the privacy audit and
PP-Mark while giving them one dissertation-level argument.  The privacy chapter
restricts what a completed execution may claim; the provenance chapter changes the
public acceptance predicate so that observable detection and verifiable process
evidence must agree.  Chapter~\ref{ch:conclusion} returns to this common argument
and states precisely what the combined evidence does and does not warrant.
